\section{Signed identity and arrival estimate}\label{app:signed}
When a queen is revealed, the uniform law weights each child by its number of completions, and the
adaptive law weights it by the predictor. Comparing the two along a history requires retaining each row
normalization and the value of the predictor at the stopping time. Let \(\mu\) be the law of the uniform
solution with an independent row order, and let \(\nu\) and \(\nu^\sharp\) be the adaptive laws for \(L\)
and \(L^\sharp\). Given row \(i\), set
\[
 \rho^\sharp(B,i)=\sum_{e\in i}\exp(L^\sharp(B+e)-L^\sharp(B)),\qquad
 A^\sharp=\sum_{t<\tau}\log\rho^\sharp(B_t,i_t).
\]
Here \(t\) counts the reveal steps in the sum, and all available children, including exit children, enter
the normalization. Let \(g^\sharp=\log Q-L^\sharp\). At arrival retain only the pair constraints of one
physical square \(C\) on the same rook mask. Write \(\gamma_C\) for the logarithm of the avoidance
probability under the rook law minus the corresponding change in the predictor. Let
\(\Gamma=\sum_C\gamma_C\) and \(g^\flat=g^\sharp-\Gamma\). The row-order factors are identical in the two
histories. The remaining probabilities telescope because the uniform law assigns to each prefix a mass
proportional to its number of completions. Inverting on \(F_b\), where this number is positive and the
history has positive mass under both laws, gives the exact identity
\begin{equation}\label{eq:signed-history}
 e^{-g^\sharp(B_n)}\nu^\sharp(F_b)
 =\mathbb E_\mu[\exp\{-A^\sharp-\Gamma-g^\flat\};F_b].
\end{equation}
Henceforth \(\mathbb E[X;E]=\mathbb E[X\mathbf1_E]\). To estimate the identity we need row consistency
and diffusion on the boards entering the normalization. We establish the latter property for the
physical child, including at first exit.
\begin{lema}[Scaling after one move]\label{lem:hijo-fisico}
Consider a regular parent of size \(m\) whose physical point \(p^{Q}\) is positive on its binary support,
has unit row and column margins, satisfies \(p^{Q}_{ic}\le C/m\), and has incidence gap at least
\(\gamma>0\). Fix a legal cell \(e\), form the rook minor \(M_e\), and forbid any subset of the other
cells on the two diagonals of \(e\). For sufficiently large \(m\), all these masks of size \(m-1\) have
scalings in a common diffusion and gap class, with constants depending only on \(C,\gamma\). In
particular this holds for the full child \(M_e^{\rm fis}\) at first exit.
\end{lema}
\begin{proof}
We delete one column in each retained row and up to one cell per diagonal. The loss in each row is at
most \(3C/m\), and the same holds per column. Let \(R\) be the restriction of \(p^Q\) after the deletions
and let \(\mathcal W\) be the side with fewer vertices of a cut in the remaining bipartite graph. In the
parent, the Rayleigh bound on the centered indicator gives capacity at least \(\gamma |\mathcal W|/2\),
since \(|\mathcal W|\le m\). The capacity remaining in \(R\) is at least \((\gamma/2-3C/m)|\mathcal W|\).
Orient the edges from rows to columns. The row deficits and the negatives of the column deficits sum to
zero, and each has absolute value at most \(3C/m\). With \(\lambda=(3C/m)/(\gamma/2-3C/m)<1\), each cut
supports the net demand with symmetric capacities \(\lambda R\). The flow criterion gives a signed
correction and a doubly stochastic matrix \((1-\lambda)R\le p'\le(1+\lambda)R\). Its conductance is at
least \((1-\lambda)(\gamma/2-3C/m)\). Cheeger's inequality gives a uniformly positive gap for the
normalized Laplacian.

It remains to transfer these bounds from a feasible table to the scaling of the binary support. At size
\(k\) let \(a\) be doubly stochastic, positive on the support, with \(a_{ic}\le C'/k\) and gap
\(\gamma'\). The entropy maximum on this support is interior by strict feasibility. Its equations give
the scaling \(D\) with \(\log(kD_{ic})=u_i-u_c\). By the margins, \(\sum a_{ic}\log(kD_{ic})=\sum
D_{ic}\log(kD_{ic})\ge0\). The positive part of the first sum is at most \(\sum a_{ic}kD_{ic}\le C'k\).
Subtracting a median of the potentials and integrating over the level cuts gives
\[
 \frac{\gamma'}2\sum_z|u_z|\le\sum_{ic}a_{ic}|u_i-u_c|\le2C'k .
\]
Fix \(T>8(C')^2/\gamma'\). Each line has at least \(k/C'\) neighbors, and in the entire graph fewer than
\(k/(2C')\) vertices satisfy \(|u_z|>T\). The row equation \(1=e^{u_i}\sum_{c\sim i} e^{-u_c}/k\) gives
\(e^{u_i}\le2C'e^T\), and the column equation gives \(e^{-u_c}\le2C'e^T\). The same equations with the
upper bounds give reciprocal lower bounds. Thus \(d_-/k\le D_{ic}\le d_+/k\) on the support, and
\(D\ge(d_-/C')a\) preserves a positive gap by Rayleigh. Apply this to \(p'\) with \(k=m-1\). For the
parent mask it applies directly to \(p^Q\), including on regular long arrival. The predictor
margins~\cite{pc} hold up to the stopping time by Corollary~4.12 and the inner class of~\cite[Def.~2.2
and §§4--5]{tc}. If the gap is stated with response weights, their comparability with \(p^Q\) on the
positive-vacancy branch gives the Rayleigh hypothesis.
\end{proof}

The coefficient \(\coefT=\mathfrak F-1/12\) in (T) depends on the marked row and column masses \(s,t\)
and the crossing weight \(z\). The integrated expressions in Appendix~\ref{app:rook} are holomorphic and
bounded on disks with class radii. At \(s=0\) or \(t=0\) the visible terms are independent of \(z\).
Differentiating with respect to \(z\) makes both faces vanish. If \(M_0\) bounds \(\coefT\) on a polydisk
of class radius \(r\) about the deletion point in \(z\), Cauchy's estimate in \(z\) bounds
\(|\partial_z\coefT|\) by \(2M_0/r\) on the smaller disk. Schwarz's lemma in the masses \(s,t\) supplies
the factors \(|s|/r,|t|/r\). Hence \(\partial_z\coefT=st\,O(1)\) uniformly, including when integrating in
\(z\) along the deletion segment. Both masses equal \(1/k\). At size \(k\) this gives \(O(k^{-3})\) for
the change in \(\coefT/k\) per forbidden cell. Removing a row and a column changes the masses by
\(O(k^{-1})\) and the inverse size by \(O(k^{-2})\), at cost \(O(k^{-2})\) in \(\coefT/k\). There are at
most \(2m\) forbidden cells. The lemma provides the same margins at every step. Summing the changes and
using the remainder \(O(m^{-2})\) of (T) at both endpoints gives, for every child of the regular parent,
\begin{equation}\label{eq:incremento-hijo-fisico}
 |g_T(M_e^{\rm fis})-g_T(M)|\le C/m^2,\qquad |g_T(M_e^{\rm fis})|\le C/(m-1).
\end{equation}


\paragraph{The sign of the exit term.}
The bound \eqref{eq:incremento-hijo-fisico} and consistency of the ordinary predictor give \(\lvert
A^\sharp\rvert\le C_b/n\) up to the stopping time, by summing \(O(m^{-2})\) over \(m>j\). On \(F_b\)
diffusion gives \(|\Gamma|\le C/j\). Set \(U_0=A^\sharp+\Gamma\). Lemma~\ref{lem:llegada-integrada}
bounds the cost of replacing \(e^{-g^\flat}\) by one in \eqref{eq:signed-history}, since \(e^{-U_0}\le
e^{C_b/n}\). The Taylor remainder for \(e^{-U_0}\) is \(O_b(n^{-2})\). The preliminary bound
\(R_Q(n)=O(n^{-1})\) of Lemma~\ref{lem:previa} and (T) give \(g^\sharp(B_n)=O(n^{-1})\), with a remainder
of the same order for its exponential. Thus, with \(\limsup_n|\epsilon_n(b)|\le C_\delta b\),
\[
 n g^\sharp(B_n)\nu^\sharp(F_b)
 =n\mathbb E_\mu[U_0;F_b]+n\{\mu(F_b^c)-\nu^\sharp(F_b^c)\}+\epsilon_n(b).
\]
The sign follows from \(\nu^\sharp(F_b)-\mu(F_b)=\mu(F_b^c)-\nu^\sharp(F_b^c)\).

The ratio of child probabilities under \(\nu^\sharp\) and \(\nu\), given the row, is \(e^{g_T(M_e^{\rm
fis})}/\mathbb E_{\nu,i}e^{g_T(M_{e'}^{\rm fis})}\). Every child of the regular parent satisfies
\(|g_T|\le C/(m-1)\), including the child at first exit. The logarithm of the ratio lies between
\(-2C/(m-1)\) and \(2C/(m-1)\), since the denominator averages over all children. The stopped product is
bounded by \(\exp(2C\sum_{m=j+1}^n(m-1)^{-1})=\exp(O_{b,\delta}(1))\). With \(X=Qe^{-L}\), the density of
\(\mu\) with respect to \(\nu\) is \(Xe^{A-R_Q(n)}\), where \(A\) sums the logarithms of the ordinary row
normalizations. On the macroscopic segment, the exit probability under \(\nu\) is exponentially small
by~\cite[Cor.~4.12]{tc}, and the cost weighted by \(e^AX\) is bounded by \(C\exp(-c_\delta n+C_\delta
n^{2/3}+O_\delta(n^{-1}))\). For each dyadic band of the stopping rule with floor \(j\), of scale \(m\ge
j/2\), the cost with weight \(1+X\) is \(C e^{-(c_*-2\beta)m}\), with \(c_*>2\beta\).
Lemma~\ref{lem:prescripcion-larga} bounds long arrival by \(O(e^{-2j})\) with that weight. On surviving
the macroscopic segment, \(|A|\le C_\delta/\lceil\delta n\rceil+C/j\) holds.
Lemma~\ref{lem:positividad-llegada} gives \(F_b^c\subset H_j^c\). The sum over bands and the macroscopic
bound control \(\nu(H_j^c)\) and \(\mu(F_b^c)\) at exponential order. The stopped product gives that
order for \(\nu^\sharp(F_b^c)\). Removing the indicator from the mean of \(A^\sharp\) then costs
\(o_b(1)\) after multiplication by \(n\). It follows that
\begin{equation}\label{eq:signed-mean}
 \limsup_n\left|n g^\sharp(B_n)-n\mathbb E_\mu A^\sharp
                 -n\mathbb E_\mu[\Gamma;F_b]\right|\le C_\delta b .
\end{equation}
\section{Summing the remainder at arrival}\label{app:arrival}
A single geometric square can contribute on the scale \(1/n\) relevant to \eqref{eq:signed-mean}. After we
subtract \(\Gamma\), each surviving support must satisfy the geometric constraints of two geometric squares.
This additional cost lets us sum the errors with a bound that decreases with \(b\). Consider first the
short historical arrival event \(H_j\) under \(\nu\). For a set of board squares \(Y\), we activate only the
diagonal constraints supported on \(Y\), keeping the entire rook mask. Let \(e(Y)\) be the logarithm
of the avoidance probability minus the change in the predictor, and let \(e_C(Y)\) be the same
expression using only the pairs of \(C\). The coefficient of a support \(U\) is \(\sum_{Y\subseteq
U}(-1)^{|U|-|Y|}e(Y)\). An isolated board square cancels when we pair subsets that differ only in that
square. If all pairs belong to one geometric square, the coefficient cancels against that of \(e_C\). If two
distinct geometric squares are involved, no \(e_C\) contributes, since the board squares would have to lie
at the four vertices of \(C\), whose only diagonal pairs are the two of \(C\). The remainder therefore
sums the original coefficients of supports that have no isolated board squares and contain at least two
distinct geometric squares. The subtracted part has degree at most four.

Let \(v\) be the number of board squares, \(h\) the total number of distinct row and column labels, \(D\)
the number of free coordinates, and \(p\) the support exponent. Prescribing the \(h\) lines gives the
factor \((j/n)^h\), together with \(\delta^{-2v}\exp(C_Lv^2/j)\) in the prescription bound. For each
type, the number of positions is of order \(n^D\). Write \(d_*=h-D\), \(q=p-D-1\). Two distinct
geometric squares impose two nonproportional linear equations on the row and column coordinates. Hence
\(d_*\ge2\), and the support bound gives \(q\ge0\). The scaled factor is
\[
 n\,n^D j^{-p}(j/n)^h=(j/n)^{d_*-1}j^{-q}
 \le (b/\delta_*)\delta_*^{d_*-1}j^{-q}
 \qquad(0<b\le\delta_*).
\]
We factor out \(b/\delta_*\) before summing. For \(v\le4\) there is an absolute bound on the number of
types. For noncollinear types the weighted count uses \(d_*-1+q\ge(v-4)/3\). Thus for \(5\le v\le
N=\lceil4\log n/\log\rho\rceil\), with \(\rho\ge2\) fixed, the bound after summing over types of that
size is
\[
 (b/\delta_*) C_L(C_1\delta^{-2})^v v^{d_{\rm pol}}\delta_*^{(v-4)/3}.
\]
Choose \(\delta_*\) so that \(C_1\delta^{-2}\delta_*^{1/3}<1\). For a collinear type with \(v\ge3\) we
have \(p=v+2,h=2v\), with at most \(4n\,n^v/v!\) supports. This gives \(C(C_2\delta^{-2})^v
v^2(j/n)^{v-2}\), and we decrease \(\delta_*\) once to sum this bound as well. Substituting atoms up
to degree \(N\) incurs an error of \(O(j^{-8})\) in value, and the global norm bounds the tail by
\(O_\rho(n^{-2})\). Both errors vanish after scaling for each fixed \(b\). Thus \(\limsup_n n\mathbb
E_\nu[|g^\flat|;H_j]\le C_\delta b\), with \(C_\delta\) independent of \(b\).

This bound controls a mean of \(|g^\flat|\) under \(\nu\). Equation \eqref{eq:signed-history} also
involves the exponential under \(\mu\). The density factor \(X\) allows us to use the same bound for
both, since at arrival
\[
 X\,|e^{-g^\flat}-1|=e^{g_T+\Gamma}|1-e^{g^\flat}|.
\]
On \(H_j\) the positive part of the diagonal error has a class bound, \(|g_T|+|\Gamma|\le C/j\) and
\(|A-R_Q(n)|=O_{b,\delta}(1/n)\). We apply \(|e^y-1|\le e^{y^+}|y|\) to the density \(Xe^{A-R_Q(n)}\).
On \(F_b\setminus H_j\) the weighted probability is \(O(e^{-2j})\). For \(0<X\le1\) we have \(X|\log
X|\le1/e\), and above one we use \(X|\log X|\le X(C_0+\beta j)\) in the outer class. The contribution
of this event to the left-hand side of \eqref{eq:arrival-estimate} before scaling is
\(O((1+j)e^{-2j})\). To obtain \eqref{eq:arrival-estimate} we first fix \(\delta\), the margins,
\(\rho\) and the class constants, then \(\delta_*\) and finally the cut-off. The threshold for \(n\)
may depend on \(b\). The constant \(C_\delta\) works for all these cut-offs, as required by the final
comparison of two cut-offs in Appendix~\ref{app:closing}.

\section{The interior mean at a positive cut-off}\label{app:interior}
At a fixed cut-off, \eqref{eq:signed-mean} requires us to sum the logarithms of the row normalizations
up to the stopping time and take the mean of the sum. For each child the first correction has size
\(1/m\). Before summing over \(m\), we must use its cancellation among the children in the same row.
For the four-family predictor, at every regular parent of size \(m\) on the auxiliary macroscopic
event down to \(b\), the insertion formula gives
\[
 L(B+e)-L(B)=\log[p_e(1+\delta_e)]+r_e^{\rm PC},\qquad
 \sum_{e\in i}p_e\delta_e=0,\quad
 |\delta_e|\le C_b/m,\quad |r_e^{\rm PC}|\le C_b/m^2 .
\]
Here the \(p_e\) are the coordinates of the four-family stationary point. Summing over the row gives
\(\log\rho(B,i)=\sum_{e\in i}p_e r_e^{\rm PC}+O_b(m^{-3})\). Replacing \(p_e(1+\delta_e)\) by \(p_e\)
costs \(O_b(m^{-3})\). Write \(V=\ell+\tfrac12\log\det S_{0,\rm gf}+3\sum p_e^2\) for the local
potential of the stationary expansion (\(\ell\) is the local term and the cofactor uses the fixed
gauge). Let \(\delta_e^0\) be its centered linear correction and \(P_e\) the second-order remainder
after extracting \(\Delta_e V=V(B+e)-V(B)\) and \(\delta_e^0\) from \(L(B+e)-L(B)-\log p_e\). Taylor
expansion gives
\[
 r_e^{\rm PC}=\Delta_e V+P_e+\tfrac12(\delta_e^0)^2
                   +(\delta_e^0-\delta_e)+O_b(m^{-3}).
\]
Although \(\delta_e-\delta_e^0=O_b(m^{-2})\) has the order we seek, both corrections have row mean
zero, so their difference vanishes exactly in that mean. The cubic error after summation is
\(O_b(n\sum_{m>j}m^{-3})=O_b(n^{-1})\). The potential contributes its endpoint values. The part we
must average over the whole interval is \(m^2[P_e+(\delta_e^0)^2/2]\). The table pairs each insertion
term with its contributions to the sixteen entries of this physical mean in the continuum.
\begin{center}\small
\begin{tabular}{ll}
\hline
\(-m^2 E_g/2\)& \(G_{\rm cof}\)\\
\(m^2 A_3\)&\(f2E+f2z+f2F\)\\
\(m^2 C^{(3)}\)&\(f2F+C3_{\rm cuad}\)\\
\(-m^2(E_Q+E_R)/2\)&\(f3F+f3c\)\\
\(m^2(-3p_e^2+p_e f_e)\)&\(f4=f4_{\rm expl}+f4_{L0}\)\\
\(m^2\sum_F p_i(1-t_{ie})w_i\)&\(f5c+f5F\)\\
\(-m^2\sum_{E'}p_iq_iw_i\)&\(f6\)\\
\(m^2(\delta_e^0)^2/2\)&\(f7\)\\
first part of \(Q_e\)&\(Q1+\Delta_1\)\\
second part of \(Q_e\), with \(-3m^2\sum_{E'}y_i^2\)&\(Q2+\Delta_2\)\\
\hline
\end{tabular}
\end{center}
\(f2F\) appears twice because the cubic term on the fiber and the pair part of \(C^{(3)}\) have the
same weight. The two parts of \(Q_e\) carry the factors \(-1/2\) and \(1/4\). An available line that
occurs in two factors counts only once, so the correction to the product is
\(\psi\psi'(U_\ell^{-1}-1)\). This gives \(\Delta_1,\Delta_2\). The twelve ordered pairs of distinct
families in the quadratic block contribute \(12/4=3\) times \(m^2\sum_{E'}y_i^2\), canceling the
indicated subtraction at each parent. These identities determine the multiplicities independently of
the numerical evaluation.

The uniform profile error in these scaled means on \([b,1]\) is \(\epsilon_n(b)\to0\). Its total cost
is at most \(\epsilon_n(b)n\sum_{m>j}m^{-2}=o_b(1)\). The predictable indicator of a regular parent is
taken before the row is chosen. Removing the exponentially rare auxiliary exits costs \(o_b(1)\),
since \(n|A|\le C_b\). On \(H_j\), let \(s=A-R_Q(n)+g_T+\Gamma\). There \(|s|=O_b(1/n)\) and
\(\int_{H_j}|d\mu-d\nu|=\mathbb E_\mu[|1-e^{-s-g^\flat}|;H_j]=O_b(1/n)\) by
\eqref{eq:arrival-estimate}. Thus we may use the adaptive law in the Riemann sum. The local potential
telescopes, with scaled endpoint contribution \(\mathcal V(b)/b-\mathcal V(1)\) and error
\(O_b(n^{-1/2})+o_b(1)\).

For the exact rook correction let \(u_e=g_T(M_e)-g_T(M)\), where \(M_e\) here denotes the mask of the
physical child. Thus \(\log\rho^\sharp-\log\rho=\log\mathbb E_{\nu,i}e^{u_e}\). Since
\(|u_e|=O(m^{-2})\), replacing the logarithm by the mean costs \(O(m^{-4})\). The weighted exit costs
and the stopped density give a stopped relative entropy of \(O_b(1)\). The conditional Pinsker
inequality and Cauchy--Schwarz bound the scaled change in the means on passing to \(\mu\) by
\(O_b(n(\sum_{m>j}m^{-4})^{1/2})=O_b(n^{-1/2})\). The increments telescope. The rook endpoints add
\((F(b)-1/12)/b-(F(1)-1/12)\) to the endpoint contribution of the local potential. Together they give
the value of \(G/x\) at \(b\) minus its value at \(1\). Writing this difference as the integral of the
derivative of \(-G/x\) gives the interior mean needed in the final comparison,
\[
 \lim_n n\mathbb E_\mu A^\sharp=\int_b^1 d(x)x^{-2}\,dx+\frac{G(b)}b-G(1)
       =\int_b^1 k(x)x^{-2}\,dx .
\]

\subsection{How to read the row expansion}
The weight of a square depends on how it enters the mean in the preceding table. For terms that
average over a center \(c\), we use the measure \(\nu(c)\,dc\) and write \(\langle\cdot\rangle_c\). A
second independent square \(i\) carries the same measure. If \(i\) lies on a line \(l\) through the
center, only one coordinate is free, either row or column, and the measure on that fiber is
\(d\omega_l(i)=\eta_l\nu(i)\,dy\). In this paragraph \(C\) denotes a kernel distinct from the column
family. Define
\begin{align*}
 A_{lh}&=(L_0-K)_{lh},\qquad B_{lh}=K_{lh}(W_{lh}-K_{lh}),\\
 C_{lh}&=-B_{lh}(W_{lh}-2K_{lh}),\qquad
 w(i)=\sum_{\{l,h\}\subset J(i),\,l\ne h}(L_0)_{lh}.
\end{align*}
Here \(W_{lh}\) equals \(W\) at the crossing, or zero if there is no crossing. We write
\(A(i,c)=\sum_{l\in J(i),\,h\in J(c)}A_{lh}\) and similarly for \(B,C\). If the first argument is a
single line, we sum only over \(h\). After summing, we set \(M_2=A^2+B,\ M_3=A^3+3AB+C\). The cubic
term distinguishes three averages according to whether we choose an independent square, a diagonal
line or a square on a fiber through the center. With their respective signs we obtain
\begin{align*}
 f2E&=\left\langle -M_3(i,c)/6+\tfrac32 W(i)M_2(i,c)\right\rangle_{i,c},\\
 f2z&=-\tfrac16\left\langle
 \int_{l\in D\cup A}\frac{u_l}{U_l}M_3(l,c)\,d\rho_l\right\rangle_c ,\\
 f2F&=\left\langle\sum_{l\in J(c)}\int_{i\in l}W(i)^2\,d\omega_l(i)\right\rangle_c .
\end{align*}
The measure \(d\rho\) is the continuum measure represented by the line weights in
Section~\ref{sec:profile-data}. Four other terms use the same averages.
\begin{align*}
 G_{\rm cof}&=-\tfrac12-3\langle W^2\rangle_c,\qquad
 f4_{\rm expl}=-9\langle W^2\rangle_c,\\
 f4_{L0}&=\langle Ww\rangle_c,\qquad f5c=-3\langle Ww\rangle_c .
\end{align*}
For \(C3_{\rm cuad}\) we consider a center and three translated squares that complete the geometric
square. We take both oriented bases (row--column and diagonal--antidiagonal) and both signs in each
direction. The weight is the product of the four \(W\) factors, \(x^{n_{RC}-3}\) and the \(U_l\)
factors of the distinct diagonal lines. Here \(n_{RC}\) counts the distinct rows and columns. We
integrate with respect to Lebesgue measure \(dr\,dc\,dh\), with \(h>0\) and the square inside the
board. The four \(W\) factors already supply the physical weight, with no additional density at the
center. In the grid implementation a side of length zero carries half weight. In \(f3F\) the quadratic
fiber term includes the variance and the covariance due to shared lines, and \(f3c\) collects the two
center contractions. In \(f5F\) we integrate over the fiber the product of \(-w(i)\) and the sum of
its line responses. In \(f7\) we square \emph{after} summing the responses on the fiber and outside
the center. Finally, \(Q1,Q2\) take the linear and quadratic trace terms with respect to the same
center, with the corrections \(\Delta_1,\Delta_2\) to the product over distinct lines given above.
These measures and multiplicities determine the interior integrand \(d(x)\). At arrival we must
average over the diagonal pairs that remain available in the rook mask.


\section{The three boundary means}\label{app:boundary} In \(\Gamma\) each square has been isolated, but
its two diagonal pairs can be available simultaneously. To preserve the sign and count their interaction
once, we start with the correction for one pair and then combine both pairs in the square. Here \(p_e\)
denotes the stationary \emph{rook} occupancy on the binary mask. For an available diagonal pair
\(\parmar=\{e,f\}\), let \(\pi^{\rm ex}_{\parmar}\) be its joint probability in a uniform rook placement.
Set \(\pi^T_{\parmar}=\exp[L_T(M_{\parmar})-L_T(M)]\), where \(M_{\parmar}\) deletes its four row and
column endpoints. With \(\mathsf G(e,f)\) the cross response of the stationary rook point and \(\chi_e\)
its single-square relative correction, define \[
 \pi_{1,\parmar}=p_ep_f(1-\mathsf G(e,f)+\chi_e+\chi_f),\qquad
 X_{\parmar}=p_ep_f(p_e+p_f)\mathsf G(e,f).
\] For a single pair the square correction is \(\log(1-\pi^{\rm
ex}_{\parmar})-\log(1-\pi_{1,\parmar})-X_{\parmar}+O(j^{-5})\). The probabilities are \(O(j^{-2})\) and
their difference is \(O(j^{-4})\). Integrating the logarithmic derivative gives a linearization error of
\(O(j^{-6})\). The full square contributes the corrections for its two pairs and an additional term
\(-3\prod_{s\in\square}p_s+O(j^{-5})\). There are at most \(j^2(j-1)\) pairs, since each square has at
most \(2(j-1)\) neighbors and each pair has two endpoints. Thus the following quantity \(\mathcal
M_n(b)\) differs from \(n\mathbb
 E_\mu[\Gamma;F_b]\) by \(O(n/j^2)=o_b(1)\). \begin{align*}
 \mathcal M_n(b)={}&-n\mathbb E_\mu[\textstyle\sum_{\parmar}(\pi^{\rm ex}_{\parmar}-\pi^T_{\parmar});F_b]\\
           &-n\mathbb E_\mu[\textstyle\sum_{\parmar}(\pi^T_{\parmar}-\pi_{1,\parmar}+X_{\parmar});F_b]
            -3n\mathbb E_\mu[\textstyle\sum_\square\prod_{s\in\square}p_s;F_b].
\end{align*} Introducing \(\pi^T_{\parmar}\) in this decomposition serves two purposes. The first mean
measures the error of the rook predictor on the pair. The second retains the insertion error for a marked
pair. To average it, we need to know which of the other two vertices remain available. Let \(u_f\) be the
final vacancies of the two diagonal families, and set \(U_f(b)=u_f+b(1-u_f)\). Let \(v_b=U_D(b)U_A(b)\),
let \(p_b\) be its doubly stochastic scaling~\cite{sk}, and put \(W_b^0=p_b/v_b\). A square with minimum
coordinates \((r,c)\) and side length \(h\) is integrated over \(\Omega=\{0<h<1,\ 0<r,c<1-h\}\). With
\(d_s=U_D(b,r-c+sh)\) and \(a_s=U_A(b,r+c+sh)\), the availability of its oriented pair contributes the
factor \(B_o\). The other two vertices have indicators \(I_l I_u,I_l I_v\). Their four joint
probabilities and the factor for the marked pair are \[
\begin{array}{c|cccc}
o&B_o&l&u&v\\\hline
 D&d_0a_0a_2&a_1&d_{-1}&d_{+1}\\
 A&a_1d_{-1}d_{+1}&d_0&a_0&a_2
\end{array}
\qquad
\begin{array}{c|c}
s&q_{o,s}\\\hline
00&1-l+l(1-u)(1-v)\\
10&lu(1-v)\\
01&l(1-u)v\\
11&luv
\end{array}
\] Each state has submeasure \(b^4 B_o q_{o,s}\,dr\,dc\,dh\). The four live labels are two rows and two
columns. A repeated diagonal vacancy retains the same indicator. Each orientation has
\(\sum_{h=1}^{n-1}(n-h)^2\) positions. The Riemann normalization contributes \(n^3\), with four factors
of \(j^{-1}\) and an outer factor of \(n\). Hence \(n^4 b^4/j^4\to1\). Inserting the two marked weights
gives \(w_{o,b}=B_oW_b^0(e)W_b^0(f)=p_b(e)p_b(f)/U_o(b)\). Even states of zero mass are retained, with no
division by \(q_{o,s}\).

For the first mean, let \(m_{f,b}\) be the mass of \(p_b\) on each line of family \(f\), with the same
fiber measure \(dr\) used for the final mass \(m_f\) in Section~\ref{sec:profile-data}. Set \[
 W_b=\sum_{f=D,A}\int\frac{m_{f,b}^2}{2U_f},\qquad
 \gamma_b=\frac{m_{D,b}}{U_D}+\frac{m_{A,b}}{U_A},\qquad
 \theta_b(r)=\int p_b(r,c)\gamma_b(r,c)\,dc .
\] \paragraph{The increment of a double minor.} Fix an admissible cut-off \(b>0\) and \(j=\lfloor
bn\rfloor\). On a regular binary mask \(M\) of size \(j\), write \(A_R,A_C\) for the uniform measures on
the live rows and columns. Let \(\mathfrak F(M)\) be the rook functional in Section~\ref{sec:torres}.
When differentiating it in its measures we retain the input mask and recompute its scaling so that both
margins remain equal to one. Write
\[
 h_R^M(i)=D_R\mathfrak F(M)[\delta_i-A_R],\qquad
 h_C^M(l)=D_C\mathfrak F(M)[\delta_l-A_C].
\]
Each influence has mean zero with respect to its own measure, by linearity. For the pair
\(\parmar=\{e,f\}\) write \(H_M(\parmar)=\sum_{z=e,f}(h_R^M(r_z)+h_C^M(c_z))\). In the minor
\(M_{\parmar}\) the measures assign mass \(1/(j-2)\) to each remaining line. In particular,
\[
 A'_R-A_R=\frac{2A_R-\delta_{r_e}-\delta_{r_f}}{j-2},\qquad
 A'_C-A_C=\frac{2A_C-\delta_{c_e}-\delta_{c_f}}{j-2}.
\]
This is an exact measure identity. Put \(\Delta_{\parmar}\mathfrak F=\mathfrak F(M_{\parmar})-\mathfrak
F(M)\). The uniform bound on the second derivative in these directions (Cauchy's estimate in the masses
of the distinguished lines) gives, for every pair in the class,
\[
 \Delta_{\parmar}\mathfrak F=-\frac{H_M(\parmar)}{j-2}+O(j^{-2}),\qquad
 j\Delta_{\parmar}\mathfrak F=-H_M(\parmar)+O(j^{-1}).
\]
Let \(D\) be the discrete doubly stochastic scaling of \(M\), and let \(w_{\parmar}=(jD_e)(jD_f)\), which
is bounded in the class. The ratio of the exact pair probability to its rook approximation and the
approximation itself satisfy
\[
 \frac{\pi_{\parmar}^{\rm ex}}{\pi_{\parmar}^T}=\exp\{g_T(M_{\parmar})-g_T(M)\},\qquad
 \pi_{\parmar}^T=j^{-2}w_{\parmar}(1+O(j^{-1})).
\]
Indeed, the first equality divides the ratios of the permanents and of their predictors. The second uses
the leading term of the two-site marginal and \(p_z/D_z=1+O(j^{-1})\). Upon subtracting
\(\coefT(M_{\parmar})/(j-2)-\coefT(M)/j\), with \(\coefT=\mathfrak F-1/12\), the two successive
increments of the remainder of \((T^+)\) are \(O(j^{-3})\). Thus,
\[
 j^2\{g_T(M_{\parmar})-g_T(M)\}
 =j\Delta_{\parmar}\mathfrak F+2\mathfrak F(M)-1/6+O(j^{-1}).
\]
Expanding the exponential in the exact ratio then gives
\begin{equation}\label{eq:pt-par-fisico}
 -j^4(\pi_{\parmar}^{\rm ex}-\pi_{\parmar}^T)
 =w_{\parmar}\{1/6-2\mathfrak F(M)-j\Delta_{\parmar}\mathfrak F\}+O(j^{-1}).
\end{equation}
The remainder is uniform over pairs. There are at most \(j^2(j-1)\) of them, and the first mean in
\(\mathcal M_n(b)\) multiplies this sum by \(n/j^4\). The total cost of this remainder and of replacing
\(-j\Delta_{\parmar}\mathfrak F\) by \(H_M(\parmar)\) is therefore
\[
 O\!\left(\frac{n}{j^4}\,j^2(j-1)\,\frac1j\right)
 =O(n/j^2)=o_b(1).
\]
 \paragraph{The density of the removed lines.} The first mean in \(\mathcal M_n(b)\) sums over the states
of both crossings and \(\sum_s q_{o,s}=1\). The pattern law and the scaling at regular arrival, developed
in Appendix~\ref{app:transfer}, give the pair measure \(w_{o,b}\,dr\,dc\,dh\) for each orientation.
Indeed, the four axial labels contribute \(b^4\), each distinct diagonal contributes a vacancy, and the
two factors \(jD_z\) converge to \(W_b^0(z)\) on the support. The outer factor \(n/j^4\) and the \(n^3\)
from the Riemann sum leave \(n^4 b^4/j^4\to1\). By symmetry between the endpoints of the pair, on each
line with coordinate \(t\) its mass integrates to \(\frac12\iint p_b(z)p_b(z')\,dr\,dr'/U_o(t)\). Summing
over lines gives \(W_b\).

If each pair contributes a test function \(a\) at each of its two rows, the symmetry of \(a(r)+a(r')\)
cancels this factor \(1/2\) when one integrates a single endpoint. In the family \(o\), integrating the
other endpoint gives \(m_{o,b}(t)/U_o(t)\). Thus,
\[
 \sum_o\int_\Omega w_{o,b}\{a(r_e)+a(r_f)\}\,dr\,dc\,dh
 =\iint a(r)p_b(r,c)\gamma_b(r,c)\,dr\,dc
 =\int a(r)\theta_b(r)\,dr .
\]
Transposition gives the same density for columns. In particular, \(\int\theta_b=2W_b\). The row measure
lost to the pairs has density \(\theta_b\), and renormalizing the surviving rows restores the mass
\(2W_b\) at constant density. Together the two operations give the direction \(\eta_b=2W_b-\theta_b\).

The centering can be verified exactly on a deterministic grid with \(m\) points per side. Let \(p_m\) be
the kernel of \(v_b\) scaled with the uniform measure \(A_m\) of mass \(1/m\) per point, and let
\(W_m,\theta_m\) be the grid sums corresponding to \(W_b,\theta_b\). On a line \(\ell\) of family \(o\)
we give each ordered pair the weight \(\omega=p_m(e)p_m(f)/(2m^3 U_o(\ell))\), including the pairs
\(e=f\) in each family. Summing over both endpoints gives \(\sum\omega=W_m\). Fixing the row \(i\) of the
first endpoint and multiplying by two gives \(\theta_m(i)/m\). With \(h_R^m,h_C^m\) denoting the centered
influences on this grid and \(H_m(\parmar)\) their sum over the four endpoints, linearity gives
\[
 \sum\omega H_m(\parmar)
 =\frac1m\sum_i\theta_m(i)(h_R^m(i)+h_C^m(i))
 =-D\mathfrak F_m[\eta_m A_m,\eta_m A_m],\qquad \eta_m=2W_m-\theta_m .
\]
The equality holds for the first variations. The pairs \(e=f\) in this quadrature have total mass
\(O_b(m^{-1})\). The physical minors always use two distinct squares.

\begin{proposicion}[First mean at fixed cut-off]\label{prop:pt-corte}
Keep \(v_b\) fixed and evaluate \(\mathfrak F\) with the two measures \((1+\varepsilon\eta_b)\,dr\),
\((1+\varepsilon\eta_b)\,dc\), rescaling \(v_b\) at each \(\varepsilon\). Denote the derivative at zero
by \(dF_2(b)\). Then
\begin{align}
 \lim_{n\to\infty}\frac n{j^4}\mathbb E_\mu
 \Big[\sum_{\parmar} w_{\parmar}\,j\Delta_{\parmar}\mathfrak F\,;F_b\Big]&=dF_2(b),
 \label{eq:pt-menores-limite}\\
 P_T(b):=\lim_{n\to\infty}-n\mathbb E_\mu
 \Big[\sum_{\parmar}(\pi_{\parmar}^{\rm ex}-\pi_{\parmar}^T)\,;F_b\Big]&=(1/6-2F(b))W_b-dF_2(b).
 \label{eq:pt-frontera-fija}
\end{align}
\end{proposicion}
\begin{proof}
To replace \(H_M(\parmar)\) by the profile influences we must also transport the derivative of the
scaling. In the series for the functional the centered kernel \(\centrado=jD-1\) has gap
\(\|\centrado\|_{2\to2}\le r_b<1\) and sections of norm at most \(C_b\). A walk of length \(k\ge2\)
between fixed endpoints is bounded by \(C_b^2r_b^{k-2}\). With \(T_R=(I-\centrado\centrado^*)^{-1}-I\),
the diagonal of \(T_R-\centrado\centrado^*\) already excludes the unique return of length two.
For columns the roles of \(\centrado\) and \(\centrado^*\) are interchanged. The truncated terms of the
functional are therefore sums over simple graphs. Differentiating a measure fixes a vertex and subtracts
its mean. Differentiating the scaling in a row \(i\) adds a path to that root, since the column response
in the centered gauge is \(-(I-\centrado^*\centrado)^{-1}\centrado_{i,\cdot}\). The choices of
differentiated vertex or edge multiply the geometric tail by a polynomial in the length, so it remains
summable.

At fixed truncation, accidental coincidences of labels cost \(O_{b,k}(1/n)\). The pattern law on live
lines, uniform in the root, takes each sum to the same graph with kernel \(p_b\). At fixed precision the
exceptional event on \(F_b\) has exponentially small mass, and the scaled weights are bounded. We take
\(n\to\infty\) and then remove the truncation. Thus in the weighted mean \(\mathfrak F(M),h_R^M,h_C^M\)
are replaced by \(F(b),h_R,h_C\), with error \(o_b(1)\). Here \(h_R,h_C\) are the continuous influences,
with mean zero. Integrating over the four endpoints gives their contribution
\[
 \int h_R\theta_b\,dr+\int h_C\theta_b\,dc
 =-D\mathfrak F[\eta_b\,dr,\eta_b\,dc]=-dF_2(b).
\]
This proves~\eqref{eq:pt-menores-limite} by \(j\Delta_{\parmar}\mathfrak F=-H_M(\parmar)+O(j^{-1})\).
Substitution into~\eqref{eq:pt-par-fisico} and use of the mass \(W_b\)
prove~\eqref{eq:pt-frontera-fija}.
\end{proof}
Identity~\eqref{eq:pt-frontera-fija} holds for the limit of physical arrival on the event \(F_b\).
If the four
axial labels are fixed and the three diagonals are required to be empty in the complete solution, the
difference between the two means is \(O(b)+o_b(1)\). Uniform positivity of \(v_b\) and continuity of the
scaling and of its first variation give \(P_T=\lim_{b\downarrow0}P_T(b)\). In this term the line measures
are varied. In \(F'(0)\), which enters the interior defect, the kernel \(v_b\) is differentiated along
the profile with both measures fixed.
 \paragraph{The third mean.} For
the third mean, the square uses six distinct diagonal labels. Two of them each occur twice among the
eight vacancy factors in \(\prod_s W_b^0(s)\). Multiplying by the six distinct vacancies and setting
\(p_0=\nu_1\) leaves exactly \[
 -S=-3\int_\Omega
       \frac{\prod_{s\in\square}\nu_1(s)}
            {u_D(r-c)\,u_A(r+c+h)}\,dr\,dc\,dh .
\] This identifies the first and third means. The second retains two pieces of mask data at the marked
square, the indicators of its two intersections. Appendix~\ref{app:weights} computes the pair error with
these two indicators fixed.

\section{The single-line functional on the marked square}\label{app:weights}
The second boundary mean compares joint insertion with its first-order approximation. Inserting \(e\)
changes the weight of \(f\) because the insertion deletes the row and column of the inserted square.
We order \(e,f\) by increasing row and retain the two crossing indicators \(X,Y\). Here \(M_e\) is the
single rook minor obtained by deleting the row and column of \(e\). Primes denote evaluation on that
minor. Let \(t=p'_f/p_f-1\), \(G_{ef}=\mathsf G(e,f)\) and \(D=\chi_e+\chi_f-G_{ef}\). The exact
product of the insertion ratios gives
\[
 e^{u}=\frac{\pi^T_{ef}}{p_ep_f}
      =(1+\chi_e)(1+\chi'_f)(1+t)\exp(r_e+r'_f).
\]
The rook increment estimate is \(r_i=-\mathsf E_i+O(j^{-3})\), also on the single minor. With
\(t=-G_{ef}+O(j^{-2})\) and \(\chi'_f-\chi_f=O(j^{-2})\), the physical remainder divided by \(p_ep_f\)
is \(\eta+O(j^{-3})\), where
\begin{align*}
 \eta={}&(t+G_{ef})+(\chi'_f-\chi_f)+\chi_e\chi_f+t(\chi_e+\chi_f)\\
        &-\mathsf E_e-\mathsf E'_f+(p_e+p_f)G_{ef}\\
     ={}&B_{\rm phys}+D^2/2+O(j^{-3}),\\
 B_{\rm phys}={}&u-D+(p_e+p_f)G_{ef}.
\end{align*}
These seven terms apply in every state and orientation. The square term \(D^2/2\) requires only the
linear term of the insertion. To identify the remainder \(B_{\rm phys}\), we must compare the action
and prefactor on the mask and its double minor. To do this within one family we remove only \(e,f\)
and retain their four lines. Both margins of square \(i\) become \(1-x_i\), and all other margins
equal one. Let \(F_-(x),\ell_-(x)\) be the stationary action and the rook correction of this punctured
mask. Set \(R(x)=F_-(x)+2\sum_i(1-x_i)\log(1-x_i)\) and \(A(x)=\ell_-(x)+\sum_i\log(1-x_i)\). Both
regularized expressions extend to \(x=(1,1)\), where the visible system is the double minor and the
total mass of its Hessian preserves the correct size constant. Let \(p=(p_e,p_f)^\top\),
\(\mathbf1=(1,1)^\top\) and \(D_p=\operatorname{diag}(p_i(1-p_i))\). For the punctured margin Hessian
\(T_-\) and the incidence \(a_i\), we have \(R''_{il}=2\delta_{il}/(1-x_i)-a_i^\top T_-^+a_l\).
Reinserting the squares adds \(\sum_i[p_i(1-p_i)/(1-2p_i)]a_i a_i^\top\). The determinant lemma on the
complement of the gauge and stationarity give
\begin{align*}
 u={}&R(\mathbf1)-R(p)-(\mathbf1-p)^\top\nabla R(p)\\
     &+A(\mathbf1)-A(p)+\tfrac12\log\det(I-R''(p)D_p).
\end{align*}
Let \(\partial=\partial_e+\partial_f\), \(H=R''(0)\) and \(\Delta C=C(\mathbf1)-C(0)\).
Differentiating the stationary equation with respect to the two weights gives the response
\((D_p^{-1}-R''(p))^{-1}\). We expand the determinant and the responses of \(\chi_e,\chi_f,G_{ef}\).
The terms linear in \(p_i H_{il}\) cancel by \(\mathbf1^\top Hp=\sum_i p_iH_{ii}+(p_e+p_f)H_{ef}\).
This leaves
\[
 B_{\rm phys}=\Delta R-\partial R(0)-\partial^2 R(0)/2+\Delta A-\partial A(0)+O_b(j^{-3}).
\]
To obtain the coefficients, we add two copies of opposite signs to each marked line. On the positive
copy \(\lambda_i=1-x_i\), and on the ordinary part and the negative copies \(\lambda_i=1\). With
auxiliary mass \(s\), the counting measure is \(M=a+s\rho\) and the margin measure is \(N=\lambda M\),
with \(\rho=0,-1,+1\) on the ordinary part, the negative copies and the positive copies, respectively.
At \(s=1/j\) the negative copies cancel the originals exactly. For the binary incidence \(Y_{il}\) and
its Sinkhorn factors~\cite{sk} \(W_{il}\), the normalized equations are
\[
 w_{il}(1-s\lambda_i\lambda_l w_{il})=W_{il}\exp(z_i+y_l),\qquad
 \int Y_{il}w_{il}\,dN_C(l)=\int Y_{il}w_{il}\,dN_R(i)=1 .
\]
Let \(R_q,A_q\) denote the coefficients of \(s^q\) in the regularized action and prefactor along
\(x_e=x_f=t\). In the action, the coefficient of \(s^0\) is affine and \(R_1\) has degree at most two
in \(t\). In the prefactor \(A_0\) is constant and \(A_1\) is affine, since at each marked endpoint
the two quadratic terms from the local factor and the margin diagonal have opposite signs and equal
magnitude. The rook gap and the kernel section bounds control the inverse of the margin Jacobian in
the gauge. The implicit equation on a fixed complex disk in \(s\), for \(|x_i|\le2\), and the Cauchy
estimates give uniform \(O_b(j^{-3})\) remainders for the regularized values and derivatives used
above. With primes denoting derivatives in \(t\), we obtain
\begin{align*}
 L_s&=A'_1(0)+\tfrac12R''_1(0),\\
 B_{2,s}&=\tfrac12\int_0^1(1-t)^2R'''_2(t)\,dt+\int_0^1(1-t)A''_2(t)\,dt .
\end{align*}
In \(L_s,B_{2,s}\) the subscript denotes the state of the crossings. The polynomials
\(R_1,R_2,A_1,A_2\) have degrees at most \(4,6,6,8\), so nine distinct nodes determine their
derivatives algebraically. Floating-point evaluation requires its own error control. To complete the
proof we must transfer \(L_s\) and \(B_{2,s}\) to the profile in each state, the first squared and the
second linearly.

\section{From marked masks to the profile}\label{app:transfer}
Two squares can depend on the survival of the same line. Treating their indicators as independent
would change the pair mean, even if the profile of each square were exact. We therefore transfer the
weighted contribution of each state with its own mass. For a fixed pattern we identify equal affine
labels and take each distinct indicator only once. With \(d_*\) free indices, an accidental
coincidence satisfies one additional nonconstant affine equation and costs \(O(n^{-1})\) after
division by \(n^{d_*}\). On the auxiliary regular trajectory, at \(x\ge b\), the hazards for surviving
axial and diagonal lines are \(1/m\) and \(\ell_f/m\), respectively. The latter differ from their
deterministic values by at most \(\eta/n\) at fixed precision. Two distinct lines are occupied
simultaneously with probability \(O_b(n^{-2})\). The Bonferroni inequalities give the product of their
one-step survival probabilities with error \(O_{b,v}(\eta/n+n^{-2})\). We multiply each current
indicator by its deterministic future survival probability up to the cut-off. One piece changes
\(O_v(n^{d_*-1})\) summands. Azuma's inequality for the centered mean, with jumps bounded by
\(O_{b,v}(1/n)\), and the regularity bound give an exponential tail at fixed precision. The union over
marked lines has polynomial size. The density ratios up to the stopping time preserve the tail,
including the bound \(\exp(O_b(n^{2/3}))\) for \(X\) when changing to \(\mu\). Thus each surviving
axial label contributes \(b\) and each distinct diagonal label contributes \(U_f(b)\) to the
normalized mean.

In particular, summing \(YW_b^0/j\) along each surviving line gives margins \(1+o_{\Pr}(1)\)
uniformly, and two surviving columns have \(\ge c_bj\) common neighbors. In logarithmic column
potentials, the map that normalizes first rows and then columns has a stochastic derivative with
entries \(\ge c'_b/j\) on the unit ball in the oscillation seminorm. Its contraction in oscillation
gives \(jD_{il}=Y_{il}W_b^0(i,l)(1+o_{\Pr}(1))\) for the Sinkhorn scaling, uniformly on the support.
When the two distinguished squares are removed, each margin loses \(O_b(1/j)\) and two columns lose at
most two common neighbors. The same proof still applies. The stationary rook occupancy satisfies
\(p_{il}/D_{il}=1+O_b(1/j)\) on the support. In the integrated patterns, the incidences from an
additional row \(y\) to the marked columns have a product law with parameters
\(v_b(y,c_0),v_b(y,c_1)\). Distinct columns give distinct diagonal labels in each family, except for
the \(O(1)\) values of \(y\) that coincide with a distinguished label. We may multiply the mean by any
of the four bounded state indicators. For the quadratic term of \(\eta\), let \(B=p_b-1\), \(G_R=BB^*,
T_R=G_R(I-G_R)^{-1}, U=T_RB\), with the transposed operators for columns. Let
\(d_R=\operatorname{diag}(T_R-G_R)\) and let \(d_C\) be its column analogue, with \(\bar d=\int
d_R=\int d_C\). Define \(\psi_R=2[(B+U)d_C-T_Rd_R]-d_R\) and its column counterpart. At the four
marked endpoints (two rows \(r_0,r_1\) and two columns \(c_0,c_1\)), the derivative identities give
\begin{align*}
 R''_1(0)/2={}&XW^0_{01}+YW^0_{10}+\sum_{\square} U-T_R(r_0,r_1)-T_C(c_0,c_1)
              -\tfrac12\sum_* d,\\
 A'_1(0)={}&-1+2\bar d+\tfrac12\sum_* \psi .
\end{align*}
The asterisk ranges over the four endpoints and the square symbol over its four board squares, with
the columns ordered so that the distinguished squares are \(00,11\). For the first identity we
differentiate the margins with respect to the occupancy parameter. On the ordinary types the mean
response solves \((I+N)y=N\alpha\), where \(N\) has off-diagonal blocks \(p_b,p_b^*\) and \(\alpha\)
inserts the two endpoints on each side. The centered blocks of \(N(I+N)^{-1}N\) are
\(T_R,-U,-U^*,T_C\). At each marked vertex a repeated incidence contributes its Bernoulli variance.
For the prefactor we differentiate the local factor, the margin diagonals, the total mass and the
centered determinant. Their sum is \(1+\langle d-1,y\rangle-\sum_*d/2\), and substituting the centered
solution and its constant mode gives the second identity. The limiting linear part \(jD\) is the sum
\(L_s\) in each submeasure. Squaring gives \(L_s^2/2\).

The pattern argument applies at any fixed length, but the second coefficient contains the inverse of
the margin matrix and hence walks of arbitrary length. We truncate its expansion in centered
alternating walks. Each walk of length \(q\ge2\) between fixed endpoints is bounded by
\(C_b^2r_b^{q-2}\), by the two section norms and the gap \(r_b<1\). At second order in the logarithm
of the determinant there are at most two insertions in each trace, and the tail is bounded by
\(C_b\sum_{q>N}(1+q)^2r_b^{\max(2q-C,0)}\). At a fixed truncation, the margin recurrences attach
rooted trees to a square or an alternating trace. Coincidences of distinct free labels cost
\(O_{b,N}(1/n)\). Each repeated indicator is retained before averaging. A trace of length two on a
free edge with conditional mean \(V\) uses \(Y^2=Y\) on the outward and return steps. The excess over
the product of the means contributes \(V(1-V)\ell^2/(2d_i d_l)\), with \(\ell\) the normalized edge
weight in the Hessian and \(d_i,d_l\) its row and column degrees. The root edges already have their
indicators fixed. Thus the ordinary variance is counted once. First we take \(n\to\infty\) at fixed
truncations, then we remove the truncations. Only then do we take \(b\downarrow0\), since the
positivity of \(v_b\) gives a uniform lower bound for \(p_b^*p_b\) near zero and dominated convergence
in each integral. With \(\mathcal I_0(C)=\sum_o\int_\Omega w_{o,0}\sum_s q_{o,s,0}C_{o,s}\), we obtain
\[
 \lim_{b\downarrow0}\limsup_n\left|
 n\mathbb E_\mu[\textstyle\sum_{\{e,f\}}p_ep_f B_{\rm phys}(e,f);F_b]
 -\mathcal I_0(B_{2,s})\right|=0. \tag{F2}
\]
Thus the signed pair contribution at the boundary is \(J=-B_2+L_2\), with \(B_2=\mathcal
I_0(B_{2,s})\) and \(L_2=-\mathcal I_0(L_s^2/2)\). In state \(11\) each orientation contributes its
single-line coefficient, and the interaction of the two orientations appears only once in \(-S\).
Adding \(J\) to the other two means gives the constant \(P_T+J-S\) that enters the cut-off comparison
once the limit in \(n\) has been taken for each fixed cut-off.

\section{Comparing two cut-offs and identifying the limit}\label{app:closing}
The preceding estimates hold for each fixed admissible cut-off. Their thresholds for \(n\) may depend
on the cut-off, so passage to the endpoint still requires justification. The quantity they approximate
on the initial board is the same for all cut-offs. Let \(a_n=nR_Q(n)\), \(a_*=F(1)-1/12\),
\(H=P_T+J-S\), and let \(b>0\) be an admissible cut-off. There are admissible cut-offs arbitrarily
close to zero. At a positive cut-off define \(I(b)=\int_b^1 k(x)x^{-2}\,dx\) and
\(\Phi(b)=a_*+I(b)+H\). The three boundary identifications give \(\eta(b)=\limsup_n|B_n(b)-H|\to0\).
The interior identity holds on \([b,1]\). Combining it with \eqref{eq:signed-mean}, the boundary
expansion with error \(O(n/j^2)\) and the limit \(n g_T(M_n)\to a_*\), we obtain
\[
 \limsup_n|a_n-\Phi(b)|\le C_\delta b+\eta(b)=:\varepsilon(b).
\]
\begin{lema}[Comparison of two cut-offs]\label{lem:dos-cortes}
Let \((a_n)\) be a real sequence and, for each admissible cut-off \(b\), let \(\Phi(b)\in\mathbb R\)
satisfy \(\limsup_n|a_n-\Phi(b)|\le\varepsilon(b)\), where \(\varepsilon(b)\to0\) as the cut-off
decreases to zero. Then \(\Phi(b)\) has a finite limit \(\Phi_0\) along the admissible cut-offs and
\(a_n\to\Phi_0\).
\end{lema}
\begin{proof}
For two fixed admissible cut-offs we compare \(\Phi(b),\Phi(b')\) with the same \(a_n\), taking \(n\)
sufficiently large for both estimates. The triangle inequality and the limit superior give
\(|\Phi(b)-\Phi(b')|\le\varepsilon(b)+\varepsilon(b')\). The Cauchy criterion holds along the
admissible cut-offs and \(\Phi(b)\to\Phi_0\in\mathbb R\). For each cut-off,
\(\limsup_n|a_n-\Phi_0|\le\varepsilon(b)+|\Phi(b)-\Phi_0|\). Since the left-hand side is independent
of \(b\), letting the cut-off decrease to zero gives \(a_n\to\Phi_0\).
\end{proof}
In particular, \(I(b)\) has a finite limit along these cut-offs. Throughout this argument we have used
the physical profile, the stopped laws and their estimates at a fixed positive cut-off.

\begin{proof}[Conclusion of the proof of Theorem~\ref{teo:principal} and Corollary~\ref{cor:derivada}]
To identify the limit with the ordinary improper integral, it remains to decide whether \(k(x)/x^2\)
can have a term proportional to \(1/x\). The zeroth-order profile identity gives \(k(0)=0\). The
projected line resolvent is analytic near zero because the diagonal Schur block has the form
\(\Lambda_u+x\mathcal N\), the vacancy has a positive minimum and the axial block has a spectral gap.
The inverse and its kernel sections remain bounded on interior complex disks. Dominated integration in
the finite contractions of the densities thus gives analyticity of \(k,G\). This regularity holds for
any value of \(k'(0)\). On a sufficiently small interval \([0,x_0]\), write \(k(x)=k'(0)x+x^2 h(x)\),
with \(h\) analytic and bounded. For \(0<b<x_0\),
\[
 I(b)=I(x_0)+k'(0)\log(x_0/b)+\int_b^{x_0}h(x)\,dx .
\]
The last integral remains bounded. Convergence of \(I(b)\) along the admissible cut-offs forces
\(k'(0)=0\). From \(k'=d'-xG''\) we deduce \(d'(0)=0\) for the physical profile. The singularity of
the integrand is then removable and the improper integral converges absolutely. Consequently, the
value of \(\Phi_0\) is
\[
 \lim_{n\to\infty}n\{\log Q(n)-\log Q_{\rm pred}(n)\}
   =F(1)-\tfrac1{12}+\int_0^1\frac{k(x)}{x^2}\,dx+P_T+J-S .
\]
\end{proof}

